\subsection{Закон всемирного тяготения}
%%%%%%%%%%%%%%%%%%%%%%%%%%%%%%%%%%%%%%%%%%%%%%%%%%%%
%%%%%%%%%%%%%%%%%%%%%%%%%%%%%%%%%%%%%%%%%%%%%%%%%%%%

 \z{Два одинаковых однородных шара, соприкасаясь, притягивают друг друга с силой $F$. Как изменится сила, если увеличить массу каждого шара в $n$ раз? Плотность материала шаров не изменяется.}
%
{Увеличится в $n^{\frac{4}{3}}$ раза}
%%%%%%%%%%%%%%%%%%%%%%%%%%%%%%%%%%%%%%%%%%%%%%%%%%%%
\z{Спутник массой $m$ движется вокруг планеты массой $M$ по круговой орбите радиуса $R$ со скоростью, вдвое превышающей первую космическую скорость для данной орбиты. Чему равна и куда направлена сила тяги, которую развивает двигатель спутника?}
%
{$F=3G\frac{Mm}{R^2}$, к центру планеты}
%%%%%%%%%%%%%%%%%%%%%%%%%%%%%%%%%%%%%%%%%%%%%%%%%%%%	
 \z[{\begin{wrapfigure}[5]{r}{0.3\textwidth}%
		\vspace{-15pt}%
		\begin{flushright}%
 			\includegraphics[width=0.3\textwidth]{4.2.3}%
		\end{flushright}%
	\end{wrapfigure}%
	}]%
	{Шар радиусом $R$ и плотностью $\rho$ имеет внутри сферическую полость радиусом $r (r < R)$, центр которой находится на расстоянии $s$ от центра шара (см. рис). С какой силой будет притягиваться к шару материальная точка массой $m$, находящаяся на расстоянии $l$ от центра шара, если треугольник $ABC$ прямоугольный с прямым углом:
		\hpu {$\angle ACB$} 
		\hpu {$\angle BAC$}}
 %
{\puan $F=\frac{4}{3}\pi G \rho m \left( \frac{R^3}{l^2}-\frac{r^3}{l^2-s^2}\right)$;~\puan $F=\frac{4}{3}\pi G \rho m \left( \frac{R^3}{l^2}-\frac{r^3}{l^2+s^2}\right)$}
%%%%%%%%%%%%%%%%%%%%%%%%%%%%%%%%%%%%%%%%%%%%%%%%%%%% 
 \z{ Космический зонд ``Шумейкер'' на некоторое время должен стать спутником астероида Эрос. По расчётам он будет обращаться вокруг астероида на высоте, составляющей $n = 1/15$ радиуса Эроса, с периодом $T = 4.5 \text{часа}$. Определите предполагаемую среднюю плотность астероида $\rho$. }
%
{$\rho =\frac{3\pi(n+1)^3}{GT^2}\approx 650 \frac{\text{кг}}{\text{м}^3}$}
%%%%%%%%%%%%%%%%%%%%%%%%%%%%%%%%%%%%%%%%%%%%%%%%%%%% 
 \z[{\begin{wrapfigure}[4]{r}{0.4\textwidth}%
		\vspace{-15pt}%
		\begin{flushright}%
 			\includegraphics[width=0.4\textwidth]{4.2.5}%
		\end{flushright}%
	\end{wrapfigure}%
	}]%
	{В открытом космосе три небольших астероида из-за гравитационного притяжения сближаются друг с другом вдоль общей прямой, неподвижной относительно звёзд. Отношение расстояний от среднего астероида до крайних остаётся равным $n = 2$ вплоть до их столкновения (см рис.). Масса левого астероида равна $m_1$, масса центрального -- $m_2$. Найдите массу $m_3$ правого астероида. }
%
{$m_3=\frac{104m_1+63m_2}{19}$}
%%%%%%%%%%%%%%%%%%%%%%%%%%%%%%%%%%%%%%%%%%%%%%%%%%%%